\let\negmedspace\undefined
\let\negthickspace\undefined
\documentclass[journal]{IEEEtran}
\usepackage[a5paper, margin=10mm, onecolumn]{geometry}
%\usepackage{lmodern} % Ensure lmodern is loaded for pdflatex
% \usepackage{tfrupee} % Include tfrupee package

\setlength{\headheight}{1cm} % Set the height of the header box
\setlength{\headsep}{0mm}     % Set the distance between the header box and the top of the text

\usepackage{gvv-book}
\usepackage{gvv}
\usepackage{cite}
\usepackage{amsmath,amssymb,amsfonts,amsthm}
\usepackage{algorithm}
\usepackage{algorithmic}
\usepackage{graphicx}
\usepackage{textcomp}
\usepackage{xcolor}
\usepackage{txfonts}
\usepackage{listings}
\usepackage{enumitem}
\usepackage{mathtools}
\usepackage{gensymb}
\usepackage{comment}
\usepackage[breaklinks=true]{hyperref}
\usepackage{tkz-euclide} 
\usepackage{listings}
% \usepackage{gvv}                                        
\def\inputGnumericTable{}                                 
\usepackage[latin1]{inputenc}                                
\usepackage{color}                                            
\usepackage{array}                                            
\usepackage{longtable}                                       
\usepackage{calc}                                             
\usepackage{multirow}                                         
\usepackage{hhline}                                           
\usepackage{ifthen}                                           
\usepackage{lscape}
\usepackage{pdfpages}
% \usepackage{algpseudocode}
\begin{document}

\bibliographystyle{IEEEtran}
\vspace{3cm}

\title{24 Hour Digital Clock}

\author{EE24BTECH11053 - S A Aravind Eswar}
% \maketitle
% \newpage
% \bigskip
{\let\newpage\relax\maketitle}

\renewcommand{\thefigure}{\theenumi}
\renewcommand{\thetable}{\theenumi}

\section{Introduction}

I have built a 24 hour clock, with 2 modes, running with help of incremental logic. 

\section{Working}

The first mode is normal function mode, where it displays currently set time(\texttt{normal mode}). In mode two, we can set time being displayed(\texttt{set mode}).To switch between the modes, we press \texttt{button 1}. 

Once you enter the \texttt{set mode}, the seconds segment will start blinking. This indicates currently seconds is selected to change. To circulate between seconds, minutes and hours segments, press \texttt{button 2}. The currently selected segment will blink. To change the time displayed the current segment, press \texttt{button 3}. For hours and minutes, the time increments, meanwhile for seconds, it resets back to \texttt{00}. The clock doesn't stop ticking in this mode.

In \texttt{normal mode}, the clock runs normally and buttons \texttt{button 2} and \texttt{button 3} are disabled.

\subsection{Hardware implementation}



\end{document}
